\section{Background and Related Work}
\label{sec:litreview}

All-sky cameras---typically consisting of a CCD or CMOS sensor behind a
fisheye lens---provide hemispherical coverage of the sky in a single
exposure.  They are deployed at observatories for cloud monitoring
\citep{yoachim_smtn005_2016}, in fireball networks for meteor orbit
determination \citep{colas_fripon_2020, gardiol_cavezzo_2021}, and
increasingly for atmospheric characterization at remote sites
\citep{hanel_measuring_2018}.  Extracting quantitative information from
these images requires solving three coupled problems: geometric
calibration of the heavily distorted fisheye projection, identification of
stars in the presence of instrumental artifacts, and photometric
measurement of those stars to derive atmospheric transmission.  We review
the state of the art in each area and identify the gaps that motivate the
present work.

% ------------------------------------------------------------------
\subsection{Astrometric Calibration}
\label{sec:litreview:astrometry}

The foundational projection model for all-sky cameras was introduced by
\citet{borovicka_astrometry_1995}, who parametrized the fisheye mapping as
a polynomial in zenith angle with explicit terms for the offset between the
optical axis and the local zenith.  This model, or variants of it, remains
the basis of most subsequent work.

\citet{barghini_astrometric_2019} revisited the Borovi\v{c}ka model in the
context of the PRISMA fireball network, identifying a fundamental source of
parameter crosstalk: the optical-axis center and the zenith projection lie
close together on the focal plane, making their separate estimation
ill-conditioned.  They introduced a ``physical coordinate'' reparametrization
that decouples these two centers, achieving $\sim$10\,arcsec random
projection indeterminacy when accumulating statistics over several months of
data.  The FRIPON network's calibration \citep{jeanne_calibration_2019}
employs a 9th-degree odd polynomial in radial distance (five free
coefficients) plus geometry and asymmetry parameters, reaching 2\,arcmin
absolute and 0.7\,arcmin internal precision.  Both methods require an
approximate initial geometry---they are refinement procedures, not blind
solvers.

More recently, \citet{yin_ali_2025} applied HEALPix segmentation of
all-sky images followed by per-tile solving with Astrometry.net
\citep{lang_astrometrynet_2010}, then stitched the per-tile solutions into
a global Generalized WCS model achieving $<$0.25\,pixel accuracy within
$70^\circ$ zenith angle.  \citet{yang_astrometry_2025} generalized
narrow-field astrometric techniques to fields exceeding $20^\circ$,
reaching $\sim$0.4\,pixel accuracy on all-sky images but similarly relying
on external plate-solving infrastructure.  \citet{tian_astrometric_2022}
trained a neural network to learn the pixel$\leftrightarrow$sky mapping,
using particle swarm optimization of the Borovi\v{c}ka model to generate
training data; the learned model achieves $\sim$1\,arcmin residuals but
requires a pre-calibrated parametric model as a training seed.
\citet{antuna-sanchez_orion_2022} developed the ORION calibration tool,
which fits three parameters (zenith pixel position, azimuth offset, and a
radial distortion curve) either automatically or with manual star
identification, achieving $\sim$9\,arcmin ($\sim$1.7\,pixel) accuracy.

A common limitation across these methods is that none provides a
fully self-contained blind solve for an uncalibrated all-sky camera from a
single frame.  Astrometry.net's geometric hashing
\citep{lang_astrometrynet_2010} is designed for approximately linear
projections; applying it to $180^\circ$ fisheye images requires
segmentation and per-tile solving \citep{yin_ali_2025}, adding pipeline
complexity and discarding the geometric constraints inherent in the global
fisheye projection.

% ------------------------------------------------------------------
\subsection{Star Detection and Matching}
\label{sec:litreview:detection}

Standard source detection algorithms (SExtractor, DAOStarFinder) were
designed for narrow-field astronomical images where the vast majority of
detected sources are real celestial objects.  In all-sky camera imagery,
the situation is inverted: dome structures, lens reflections, internal
baffling features, and horizon obstructions routinely dominate the
detection list.  \citet{zhi_cloud_2024} used SExtractor for star
extraction in their cloud reconstruction pipeline, achieving 0.87\,pixel
whole-sky calibration accuracy, but did not address the artifact
contamination problem---their camera installation apparently lacked
significant dome or obstruction features.

The PRISMA pipeline \citep{barghini_astrometric_2019} mitigates this
by starting from the brightest detected sources (which are more likely to
be real stars) and iteratively extending to fainter magnitudes as the model
improves.  This top-down approach helps but still requires a detection
step, and the brightest ``sources'' in many all-sky installations are in
fact dome glow or lens reflections.

An alternative paradigm is \emph{forced} or \emph{guided} photometry, in
which catalog star positions are projected into pixel coordinates using a
known (or approximate) camera model, and flux is measured at those
predicted positions without running general source detection.  The
LSST all-sky camera team explored this approach
\citep{yoachim_smtn005_2016} but reported that their camera lacked
sufficient sensitivity to detect enough stars at adequate signal-to-noise
for full transparency mapping.  \citet{roman_aod_2025} measured flux at
the positions of 56 pre-selected reference stars across eight all-sky
cameras deployed at nine sites, demonstrating the viability of forced
photometry for aerosol optical depth retrieval but limiting spatial
resolution to $\sim$56 independent measurements per hemisphere.

% ------------------------------------------------------------------
\subsection{Cloud Detection}
\label{sec:litreview:clouds}

Cloud detection from all-sky cameras broadly falls into three categories.

\textbf{Threshold and differencing methods} detect clouds by examining
temporal changes or spatial uniformity.  \citet{yoachim_smtn005_2016}
differenced consecutive LSST all-sky frames, applied Gaussian smoothing,
and flagged pixels exceeding $3\sigma$ above the RMS.  This detects
\emph{changes} in cloud cover but cannot identify static overcast
conditions or measure cloud opacity.  \citet{buntin_cloud_2025} extended
this approach with Otsu thresholding on difference images and
morphological cleanup, achieving $\sim$1\% false positive/negative rates
on 636\,000 images from the Liverpool Telescope's SkyCam and demonstrating
cloud tracking and prediction up to 15 minutes ahead.

\textbf{Machine learning classifiers} treat cloud detection as an image
classification or segmentation problem.  \citet{mommert_cloudynight_2020}
trained ResNet and LightGBM models on $\sim$2000 manually labeled images
from the Lowell Discovery Telescope all-sky camera, with LightGBM achieving
95\% accuracy on extracted features.  \citet{li_classification_2022} used
random forests on cloud area, weight, and dispersion features, reaching
96.6\% accuracy.  The LenghuSky-8 dataset \citep{rui_lenghusky_2025}
provides 429\,000 frames with three-class cloud masks (sky, cloud,
contamination) and per-pixel alt-az calibration from an 8-year campaign at
the Lenghu site; a DINOv2 linear probe achieves 93.3\% segmentation
accuracy.  These methods provide binary or categorical cloud masks but do
not measure quantitative transmission or opacity.

\textbf{Star photometry methods} use the measured brightness of detected
stars, compared to catalog values, as a probe of atmospheric transparency.
\citet{zhi_cloud_2024} developed seven cloud-layer reconstruction methods
(meshing plus machine learning interpolation) from SExtractor photometry,
achieving 1.8$^\circ$ reconstruction accuracy.  The CTA all-sky camera
system counts detected versus expected stars per exposure as a binary cloud
indicator.  \citet{bortolotto_extinction_2014} and
\citet{anghel_photometric_2019} used all-sky camera photometry to derive
extinction coefficients via the Bouguer method, measuring global (not
spatially resolved) atmospheric properties.

% ------------------------------------------------------------------
\subsection{Atmospheric Transmission Measurement}
\label{sec:litreview:transmission}

Quantitative, spatially resolved atmospheric transmission from all-sky
cameras remains largely an open problem.  The traditional Bouguer--Langley
technique \citep{bortolotto_extinction_2014, anghel_photometric_2019}
tracks individual stars across airmasses through a night, deriving a single
extinction coefficient that assumes a horizontally uniform atmosphere.
\citet{roman_aod_2025} extended star photometry to aerosol optical depth
retrieval across a network of all-sky cameras, validating against lunar
photometers (correlation $>$0.90), but their method produces a bulk AOD per
camera rather than a spatially resolved map.

\citet{zhi_cloud_2024} produced the most spatially detailed cloud
reconstruction to date, using SExtractor-detected stars as probes of local
transmission and interpolating via meshing and machine learning.  However,
their approach inherits the limitations of source detection (artifact
contamination) and achieves relatively coarse ($1.8^\circ$) angular
resolution.  No published method addresses the problem of \emph{unmatched}
catalog stars---sky regions where no star is detected at all---as
quantitative zero-transmission anchors in the interpolation.

% ------------------------------------------------------------------
\subsection{Motivation for This Work}
\label{sec:litreview:motivation}

The gaps identified above motivate the present system.  First, no existing
method provides a fully blind, self-contained astrometric solve for an
uncalibrated all-sky fisheye camera from a single frame---all published
approaches require either external solving infrastructure, approximate
initial geometry, or multi-frame accumulation.  Second, source detection in
all-sky imagery is fundamentally compromised by dome and obstruction
artifacts, yet most photometric pipelines rely on general-purpose detection
algorithms without addressing this.  Third, spatially resolved atmospheric
transmission maps remain rare and coarse, with no method exploiting
unmatched-star probes or local-background photometry to handle the bright
foreground contamination common in dome-enclosed cameras.

We present \textsc{AllClear}, an open-source pipeline that addresses these
three challenges through: (1) a hypothesis-and-verify blind solver that
determines camera geometry from scratch without external tools; (2)
detection-free guided matching that projects catalog positions into pixel
space and searches for point-source peaks, bypassing artifact-dominated
detection lists; and (3) a transmission measurement system that combines
local-background-subtracted photometry, absolute calibration via a stored
photometric zeropoint, and zero-transmission probing at unmatched catalog
positions to produce continuous all-sky transmission maps.
